\section{Introduction}
\label{sec:intro}
Consider two npm workspace packages: a shared asset library and a UI package
that imports it, followed by two Docker images. The dependent UI package
consumes its provider's local build output, and the first image packages the
built UI. The second image uses the first as its base; when built by an isolated
registry-based builder, it must pull the new base version from a registry.
Building all packages before publishing any of them is not a valid schedule
for this release. Serializing every package through both stages is sufficient
but also forces local-artifact consumers to wait for uploads they do not need.

Interruption makes the problem harder. If both npm packages and the base image have
published but the final image fails, a repeated command should finish the
remaining release. It should neither forget that image because its providers
have already released nor assign a different version solely because the first
attempt made progress. A release engine must therefore answer three questions:
which versions are owed, which stages may run, and which obligations remain
after a failure. These questions interact, but require different state.

Release engineering spans more than compilation~\cite{adams}. A heterogeneous
release writes to independently administered package registries, container
registries, and deployment services. They provide no shared atomic commit.
The resulting execution resembles a saga~\cite{sagas}: individual steps commit
separately, and recovery must account for the work that already succeeded.
For published artifacts, deletion may be unavailable or operationally
unacceptable. Our protocol therefore uses forward recovery. It completes or
reconciles outstanding work rather than assuming that every external effect
has a compensating action.

Workflow engines can express a build/publish DAG. Existing release tools also
compute versions and coordinate publication~\cite{lerna,nx,changesets,
releaseplease,semanticrelease}. The question studied here is how to compose
stage ordering with a release ledger so that successful providers do not
erase their failed consumers' obligations. The key rule is consumer-based
admission: a propagated unit of intent remains pending according to the
consumer's own release window, even when the originating package has already
recorded its release. A second distinction separates intent accumulated over
a prerelease train from fresh work not yet delivered.

We present \tool{}, an implementation of this protocol, and \lang{}, its
release-intent language. These names are pseudonyms for anonymous review;
reported version numbers identify the evaluated artifacts. The implementation
uses Git history and tags as its release ledger and ordinary commands for
build and publish stages. It requires no additional persistent release-state
database. That claim concerns orchestration state, not the storage used by
registries, build tools, or deployed applications.

A package is a versioned release unit: software, specifications, and
documentation use the same intent, dependency, command, and completion
interfaces. The protocol accommodates different build tools and registries;
its guarantees require conforming integrations, not a particular language.
The same graph also gives provisioned CI a stable entry point for affected
tests and releases as packages evolve; configuration still records each new
package, and new credentials or runners can require workflow changes.

The paper makes three contributions:
\begin{itemize}
\item An ecosystem-independent, saga-inspired release protocol combining
local-artifact and published-artifact dependencies with explicit release
intent and per-package completion records
(Section~\ref{sec:system}).
\item A self-contained model with deterministic planning and eight recovery,
discharge, and scope/channel guarantees under stated assumptions
(Sections~\ref{sec:model}--\ref{sec:guarantees}). Counterexamples explain why
consumer windows and stable retry predicates matter.
\item An evaluation of a released binary using differential planning,
targeted vectors, injected failures, historical replay, and scale measurements
(Section~\ref{sec:eval}), with the model-to-implementation gap made explicit.
\end{itemize}

The mathematical contribution is the admission and discharge discipline,
not a new topological-sorting or Git-storage primitive.
